% ---------------------------------------------------------------------------
% Discussion
%
%   - where each paradigm wins (work item 4). H1-H6: HL at low-to-mid budget
%     when rules are complex and replays interpretable; RL asymptotically when
%     rules are simple, the state space vast, and simulation cheap.
%   - the information-channel asymmetry, stated as a finding rather than as a
%     caveat: rules and replays are a wider channel than a scalar reward, and
%     that is what the budget-conditional result is evidence for. State the
%     scope with it, in the same breath rather than in a later caveat: the
%     channel is what earns the efficiency, so the finding is not that symbolic
%     policies are informationally privileged. A memory-based agent with no
%     program gets a comparable gain (echo2025hindsight). What the program buys
%     is the other bound. Written the loose way, this bullet becomes the
%     paper's most quotable overclaim and the abstract no longer supports it.
%   - the two bounds as one claim, which is the paper's thesis and should be
%     stated here in one paragraph: HL succeeds while K(pi_good) stays under the
%     agent's maintenance capacity. Both the win and the plateau follow from that
%     one inequality, and the reader should leave with it rather than with two
%     separate findings.
%   - the consequence worth drawing out: the ceiling is a property of the
%     maintainer, not of the paradigm, so it moves with model capability. State
%     what that does and does not license. It does not license extrapolating to
%     future models -- we measured one agent generation. It does mean the
%     asymptotic argument against symbolic methods was never about symbolic
%     methods. Resist the stronger version of this; it is the most tempting
%     overclaim available and the easiest to attack.
%   - why language-scale tasks fall outside, in one sentence, from the same
%     inequality: not because they are unstructured, but because no policy of a
%     maintainable description length is any good at them.
%   - what the winning human programs actually contain, and what that says about
%     representation (RQ3).
%   - Limitations, as its own labelled subsection. Two objections are raised in
%     the introduction and have to be answered in full here, not merely repeated.
%
%     Pretraining. The efficiency comparison is between an agent carrying a
%     pretraining corpus and RL starting cold, and no denominator makes that
%     symmetric. What is defensible: environment steps are the scarce resource
%     in a new domain, and the games are the venue's own, so what pretraining
%     plausibly supplies is general coding competence rather than strategy for
%     these games. Word it as plausibility plus a pointer to the Section 4
%     contamination results -- NOT as a consequence of the rules being original,
%     which does not follow (the archive is public) and which the introduction
%     and Section 2 both stopped asserting. This bullet was the last place the
%     provenance version survived; if the check comes back showing familiarity
%     with a game, the honest form here is the conditional one, reporting that
%     game separately rather than resting the limitation on an assumption the
%     data contradicts.
%     Report the figures as conditional on a general-purpose model existing.
%     Do not attempt to amortize pretraining into a per-task number -- any such
%     figure is arbitrary and a referee will say so. If a pretraining-inclusive
%     comparison is wanted, the honest form is a sensitivity statement: how large
%     would the amortized cost have to be to erase the gap?
%
%     The selection effect: the archive's RL scarcity is confounded with student cost
%     structure -- no GPUs, no time, no RL expertise -- so it is evidence about
%     what wins under contest constraints, not about what RL can reach. Say this
%     plainly and cite the budget ledger. Burying it here unlabelled reads as
%     evasion; a reviewer who has to find it themselves will not be generous.
%     Also: single archive, single institution, agent scaffolding held fixed, and
%     Elo anchored across editions whose *human* population changes composition
%     year to year -- a real limitation, and a different pool from HL's frozen
%     test pool, which does not drift. Keep the two distinct wherever both appear.
%   - negative results are in scope and belong here.
% ---------------------------------------------------------------------------

% PRACTICAL GUIDELINES ARE NOW PROMISED HERE, by contribution 3 of the
% introduction as of 2026-08-18 ("From that we derive practical guidelines for
% applying HL -- when to expect it to pay, and when a learned method is the better
% spend"), which forward-references this section. Nothing in the paper delivers
% them yet, and a contribution the paper does not keep is worse than one it never
% claimed. What the guidelines have to be derived FROM, so they are results rather
% than opinions:
%   - the K(pi_good) inequality: expect HL to pay while the winning policy's
%     description length sits below the agent's maintenance capacity. Since
%     \val{maintcap} is measured and roughly game-independent, that is an
%     up-front test a practitioner can apply to a new game by compressing a
%     reference implementation (\val{kproxy}) -- state it that way, as a check
%     someone can run, not as a restatement of the inequality.
%   - the crossover: below \val{crossover} environment steps HL is the better
%     spend, above it model-based RL is. Attach the denominators, because HL costs
%     more \val{loseaxis} and a guideline that hides that is advertising.
%   - the plateau: near \val{plateau} revisions the loop stops paying, so budget
%     for that rather than for indefinite improvement.
% Each guideline must name the result it comes from. If a number is unmeasured the
% guideline does not get stated -- see the numbers guard in AGENTS.md.

% SCOPE SUBSECTION, added 2026-08-19 at the user's direction ("我们在衡量的时候……你要
% 强调并加上一个限定，说明是在这种博弈的情况下"). Section 1 and Section 2 both forward-
% reference this text for what fails outside the class, so it is owed here.
%
% IT CONTAINS NO RESULT OF OURS, deliberately -- it is an argument about the four
% conditions in Section 2 and what depends on each, all of which is established before
% any run. That keeps it writable now, under the numbers guard. Do NOT add a measured
% claim here later; if the scope argument acquires evidence, the evidence belongs in
% Section 5 and this text cites it.
%
% ONE THING IT MUST NOT BECOME: a general "limitations" list. Threats to validity
% (single archive, single institution, fixed scaffolding, Elo drift across editions)
% are a separate subsection per the header notes above, and merging them would bury the
% scope statement, which is the one a referee checks the title against.
\section{Discussion}
\label{sec:discussion}

\subsection{What the scope excludes}
\label{sec:scope-limits}

The four conditions of Section~\ref{sec:preliminary} are the loop's requirements, not
descriptive detail, and each names a way a task outside the class defeats it. When rules
are not fixed, a revision that was correct becomes wrong for reasons the loop cannot
distinguish from its own error, and the accumulated program stops being an asset. When
legality is not machine-checkable, a failure cannot be localised to a decision, and the
agent's read of its own loss degrades to a guess --- the bound
Section~\ref{sec:hl-bottleneck} identifies. When the verdict comes from a model or a
person rather than a referee, an improvement claim is no longer falsifiable at the
granularity a revision operates on, and the loop can optimise the judge. When no
participant lies outside the policy's control, the signal is exhaustible: a fixed task is
solved once and then supplies nothing further, which is precisely the condition under
which continued maintenance has nothing to continue on.

These are not hypothetical exclusions. The suite itself contains one task that breaks two
of the conditions --- DeepClue is single-player and its referee calls a language model to
adjudicate --- which is why Section~\ref{sec:benchmark} reports it apart from the main
claim rather than dropping it: it is the boundary case, and it is more informative
standing at the boundary than absent. Much of the work an agent is actually asked to do
sits further outside. Open-ended software maintenance has no referee, specifications that
move, and no opponent; a task graded by human preference has a judge that can be
optimised. Nothing measured here predicts what happens there, and the ceiling this paper
reports should not be read as a ceiling on policy-level continual learning in general.
The narrower reading is also the more useful one: because the class is favourable, a
limit found inside it is evidence about the agent rather than about the task, which is
what makes the maintenance-capacity account testable at all.
